\begin{enumerate}
%1
\item \textbf{การกระจัด ความเร็ว และความเร่ง}

ลูกบอลเคลื่อนที่แบบโปรเจคไทล์ โดยมีสมการความสูงจากพื้นเป็น $s(t) = -10t^2 + 40t + 50$
\begin{tasks}[after-item-skip = 40pt, after-skip = 40pt](1)
\task จงหาสมการความเร็ว $v(t)$ และความเร่ง $a(t)$
\task ตอนเริ่มต้น ลูกบอลอยู่สูงจากพื้นเท่าใด
\task ลูกบอลมีความเร็วเป็น 0 ($v(t) = 0$) เมื่อใด
\task ลูกบอลตกพื้น ($s(t) = 0$) เมื่อใด
\end{tasks}


\item  เครื่องบินเคลื่อนที่เป็นเส้นตรง โดยมีสมการการกระจัดเป็น $s(t) = 2t^3 - 12t + 18$
\begin{tasks}[after-item-skip = 40pt, after-skip = 40pt](1)
\task จงหาสมการความเร็ว $v(t)$ และความเร่ง $a(t)$
\task จงหาค่าของ $v(1)$ และ $a(1)$
\task จงหาว่าเครื่องบินเปลี่ยนทิศทางเมื่อเวลา $t$ เท่าใด
\task จงหาว่าเครื่องบินกลับมาที่จุดเริ่มต้นเมื่อเวลา $t$ ใด
\end{tasks}


\newpage 

\item \textbf{อัตราสัมพัทธ์}
\begin{tasks}[after-item-skip = 150pt, after-skip = 150pt](1)
\task ถ้าแผ่นเหล็กวงกลมขยายขึ้น โดยมีอัตราการขยายของรัศมีเป็น 2 เซนติเมตรต่อวินาที จงหาอัตราการเปลี่ยนแปลงของพื้นที่เมื่อแผ่นเหล็กมีรัศมี 5 เซนติเมตร
\task ถ้าสูบลมเข้าไปในลูกบาศก์ยางด้วยอัตราเร็วคงที่ 10 ลูกบาศก์เซนติเมตรต่อวินาที จงหาอัตราการเปลี่ยนแปลงของความยาวด้านที่เพิ่มขึ้น เมื่อลูกบาศก์มีปริมาตร 100 ลูกบาศก์เซนติเมตร
\end{tasks}

%2
\item \textbf{การหาค่าประมาณ}

กำหนดฟังก์ชัน $f(x)$ ที่หาอนุพันธ์ได้ที่จุด $x_0$ ฟังก์ชันการประมาณเชิงเส้นคือ
\[ L(x) = \rule{100pt}{0.5pt} \]

จงใช้ฟังก์ชันการประมาณเชิงเส้น เพื่อประมาณค่าต่อไปนี้
\begin{tasks}(3)
\task $(100.5)^2$ 
\task $\sqrt[3]{1001}$
\task $\frac{1}{2001}$
\end{tasks}

\end{enumerate}